\documentclass[10pt,a4paper,twoside]{article}

\usepackage[left=22mm,right=22mm,top=15mm,bottom=23mm,includehead,headheight=11pt,headsep=6mm,footskip=12mm]{geometry}
\usepackage{fontspec}
\setmainfont{TeX Gyre Termes}
\setsansfont{Latin Modern Sans}
\setmonofont{Latin Modern Mono}
\usepackage{amsmath,amssymb,booktabs,caption,enumitem,fancyhdr,graphicx,hyperref,microtype,natbib,titlesec,xcolor}
\usepackage[most]{tcolorbox}

% Zhejiang University's Qiushi Blue (Truth-Seeking Blue):
% CMYK 100, 70, 0, 25; RGB 0, 63, 136.
\definecolor{QiushiBlue}{RGB}{0,63,136}
\colorlet{ReportBlue}{QiushiBlue}
\definecolor{TextGray}{HTML}{555B66}

\hypersetup{
  colorlinks=true,
  linkcolor=ReportBlue,
  citecolor=ReportBlue,
  urlcolor=ReportBlue,
  pdfauthor={First Author, Second Author},
  pdftitle={Adaptive Spatial Reasoning with Structured Visual Tokens}
}

% Choose `report` or `paper`; the title label, author block, and running header
% adapt to the selected mode.
\newcommand{\ReportMode}{report}
% In report mode, choose `front` or `back` for the author information.
\newcommand{\ReportAuthorPlacement}{back}

% Replace these fields with the report's metadata.
\newcommand{\ReportTitle}{Adaptive Spatial Reasoning with Structured Visual Tokens}
\newcommand{\ReportShortTitle}{Adaptive Spatial Reasoning}
\newcommand{\ReportShortAuthors}{First Author et al.}
\newcommand{\ReportAuthors}{First Author\textsuperscript{1,*}, Second Author\textsuperscript{1,2,\ensuremath{\dagger}}, Third Author\textsuperscript{2}}
\newcommand{\ReportLabName}{ReLER Lab}
\newcommand{\ReportInstitution}{\ReportLabName, Zhejiang University}
\newcommand{\ReportAffiliations}{\textsuperscript{1}Zhejiang University\hspace{1em}\textsuperscript{2}Research Institute}
\newcommand{\ReportAuthorNotes}{\textsuperscript{*}Equal contribution \quad \textsuperscript{\ensuremath{\dagger}}Corresponding author}
\newcommand{\ReportNumber}{ZJU-TR-2026-01}
\newcommand{\ReportDate}{September 27, 2026}
\newcommand{\ReportContact}{Second Author at \href{mailto:author@example.edu}{author@example.edu}}
\newcommand{\ReportKeywords}{spatial reasoning, visual generation, structured representations}
\newcommand{\ZJULogo}{assets/zju-horizontal-blue.png}
\newcommand{\LabLogo}{assets/lab-logo.png}

\newif\ifpaperreport
\newif\iffrontauthorblock
\newif\ifbackauthorlist
\newcommand{\PaperReportMode}{paper}
\newcommand{\FrontAuthorPlacement}{front}
\newcommand{\BackAuthorPlacement}{back}
\newcommand{\ReportHeaderLabLogo}{}
\ifx\ReportMode\PaperReportMode
  \paperreporttrue
  \frontauthorblocktrue
  \backauthorlistfalse
  \newcommand{\ReportTypeLabel}{RESEARCH PAPER}
  \newcommand{\ReportOddLeftHeader}{\ReportShortTitle}
  \newcommand{\ReportOddRightHeader}{}
  \newcommand{\ReportEvenLeftHeader}{}
  \newcommand{\ReportEvenRightHeader}{\ReportShortAuthors}
  \newcommand{\ReportIdentifierLine}{}
\else
  \paperreportfalse
  \frontauthorblockfalse
  \backauthorlistfalse
  \ifx\ReportAuthorPlacement\FrontAuthorPlacement
    \frontauthorblocktrue
  \fi
  \ifx\ReportAuthorPlacement\BackAuthorPlacement
    \backauthorlisttrue
  \fi
  \newcommand{\ReportTypeLabel}{TECHNICAL REPORT}
  \renewcommand{\ReportHeaderLabLogo}{\hspace{6mm}\raisebox{2mm}{\includegraphics[height=5.5mm]{\LabLogo}}}
  \newcommand{\ReportOddLeftHeader}{}
  \newcommand{\ReportOddRightHeader}{\ReportShortTitle}
  \newcommand{\ReportEvenLeftHeader}{}
  \newcommand{\ReportEvenRightHeader}{\ReportShortTitle}
  \newcommand{\ReportIdentifierLine}{\quad\ReportNumber}
\fi

\titleformat{\section}
  {\rmfamily\bfseries\fontsize{12}{14}\selectfont\color{ReportBlue}}
  {\thesection}{0.65em}{}
\titleformat{\subsection}
  {\rmfamily\bfseries\fontsize{11}{13}\selectfont\color{ReportBlue}}
  {\thesubsection}{0.65em}{}
\titlespacing*{\section}{0pt}{1.8ex plus 0.5ex minus 0.2ex}{0.7ex}
\titlespacing*{\subsection}{0pt}{1.3ex plus 0.4ex minus 0.2ex}{0.5ex}
\captionsetup{font=small,labelfont=bf,skip=5pt}
\setlist{itemsep=2pt,topsep=4pt}
% Keep the English body close to the compact 10pt/11.5pt report rhythm.
\linespread{0.95}
\setlength{\parindent}{0pt}
\setlength{\parskip}{3pt}
\setlength{\emergencystretch}{2em}

\newcommand{\makezjutitle}{%
  \noindent
  \raisebox{-0.5mm}{\includegraphics[height=10.5mm]{\ZJULogo}}%
  \ReportHeaderLabLogo
  \hfill
  {\sffamily\small\color{TextGray}\ReportTypeLabel\ReportIdentifierLine}\par
  \vspace{5pt}\noindent\textcolor{QiushiBlue}{\rule{\linewidth}{0.8pt}}\par
  \vspace{13pt}
  {\centering\rmfamily\bfseries\fontsize{18}{21}\selectfont\ReportTitle\par}
  \vspace{10pt}\noindent\textcolor{QiushiBlue}{\rule{\linewidth}{0.8pt}}\par
  \vspace{9pt}
  \ifpaperreport
  {\centering\rmfamily\bfseries\fontsize{12}{14}\selectfont\ReportAuthors\par}
  \vspace{3pt}
  {\centering\fontsize{10}{12}\selectfont\ReportAffiliations\par}
  \vspace{3pt}
  {\centering\includegraphics[height=5.4mm]{\LabLogo}\par}
  \vspace{2pt}
  {\centering\itshape\fontsize{9}{11}\selectfont\color{TextGray}\ReportAuthorNotes\par}
  \vspace{10pt}
  \else
    \iffrontauthorblock
    {\centering\rmfamily\bfseries\fontsize{12}{14}\selectfont\ReportAuthors\par}
    \vspace{3pt}
    {\centering\rmfamily\fontsize{10}{12}\selectfont\ReportInstitution\par}
    \vspace{2pt}
    {\centering\itshape\fontsize{9}{11}\selectfont\color{TextGray}\ReportAuthorNotes\par}
    \vspace{10pt}
    \else
    {\centering\rmfamily\bfseries\fontsize{11}{13}\selectfont\ReportInstitution\par}
    \vspace{10pt}
    \fi
  \fi
  \begin{tcolorbox}[
    enhanced,
    colback=white,
    colframe=QiushiBlue,
    boxrule=0.8pt,
    arc=2mm,
    left=8mm,right=8mm,top=5mm,bottom=5mm,
    before skip=0pt,after skip=0pt
  ]
  % Keep every element of the abstract block in the report's serif face.
  {\rmfamily
  {\centering\rmfamily\bfseries\fontsize{12}{14}\selectfont\color{ReportBlue}Abstract\par}
  \vspace{3pt}
  {\small
    Generative visual models can follow broad textual descriptions while still
    missing exact spatial constraints. We study a structured representation
    that separates object identity from position, scale, and orientation, then
    use it to guide generation and revision. The method is designed to make
    spatial requirements inspectable before rendering and to support targeted
    corrections without regenerating the entire scene. This sample document
    demonstrates the report layout; all names, claims, and numerical values are
    placeholders and should be replaced with verified project material.\par}
  \vspace{4pt}
  {\footnotesize\textbf{Date:} \ReportDate\par
  \iffrontauthorblock
  \textbf{Correspondence:} \ReportContact\par
  \fi
  \textbf{Keywords:} \ReportKeywords\par}
  }
  \end{tcolorbox}
  \vspace{13pt}
  \thispagestyle{plain}
}

\pagestyle{fancy}
\fancyhf{}
\fancyhead[LO]{\footnotesize\rmfamily\color{TextGray}\ReportOddLeftHeader}
\fancyhead[RO]{\footnotesize\rmfamily\color{TextGray}\ReportOddRightHeader}
\fancyhead[LE]{\footnotesize\rmfamily\color{TextGray}\ReportEvenLeftHeader}
\fancyhead[RE]{\footnotesize\rmfamily\color{TextGray}\ReportEvenRightHeader}
\fancyfoot[C]{\footnotesize\rmfamily\color{TextGray}\thepage}
\renewcommand{\headrulewidth}{0pt}
\renewcommand{\footrulewidth}{0pt}
\fancypagestyle{plain}{%
  \fancyhf{}
  \fancyfoot[C]{\footnotesize\rmfamily\color{TextGray}\thepage}
  \renewcommand{\headrulewidth}{0pt}
  \renewcommand{\footrulewidth}{0pt}
}

\begin{document}

\makezjutitle

\section{Introduction}

Text-to-image systems have improved substantially, but precise spatial control
remains difficult. Large pretrained models provide a powerful basis for visual
generation \citep{radford2021clip}, yet a prompt can name the right objects and
still produce the wrong arrangement, scale, or relative position. This gap limits use cases where
the visual output must satisfy a layout or scene specification.

We represent a scene as a set of object records. Each record contains a label
and a normalized bounding box, so users can inspect and edit spatial intent
before image generation. The representation is deliberately simple: it can be
created from language, edited directly, and translated into constraints for a
downstream generator.

\section{Method}

For an image of width $W$ and height $H$, an object $i$ is represented by
$o_i=(c_i,x_i,y_i,w_i,h_i)$, where $c_i$ is its category and the remaining
values are normalized coordinates. The center of the object is
$(x_iW,y_iH)$ and its size is $(w_iW,h_iH)$. A scene is the collection
$\mathcal{S}=\{o_i\}_{i=1}^{N}$.

An instruction such as ``move the cup to the left of the plate'' can be encoded
as an ordering constraint:
\begin{equation}
  x_{\mathrm{cup}} + \frac{w_{\mathrm{cup}}}{2}
  < x_{\mathrm{plate}} - \frac{w_{\mathrm{plate}}}{2}.
  \label{eq:left-of}
\end{equation}
The layout can then be checked against explicit constraints before it is sent
to an image generator.

\subsection{A small example}

Table~\ref{tab:layout} shows example object records. The values are illustrative
and are included only to demonstrate table styling.

\begin{table}[htbp]
  \centering
  \caption{Example normalized object layout.}
  \label{tab:layout}
  \begin{tabular}{lrrrr}
    \toprule
    Object & $x$ & $y$ & $w$ & $h$ \\
    \midrule
    Cup   & 0.28 & 0.62 & 0.12 & 0.20 \\
    Plate & 0.66 & 0.62 & 0.26 & 0.18 \\
    \bottomrule
  \end{tabular}
\end{table}

\section{Evaluation}

A useful evaluation should measure both visual quality and constraint
satisfaction. For example, a benchmark can report the fraction of generated
scenes that satisfy all specified relations, alongside human preference or a
perceptual quality metric. Reporting only an aggregate image score can obscure
whether the system followed the requested layout.

\section{Limitations}

Bounding boxes do not fully describe occlusion, articulated pose, or complex
three-dimensional relations. Some instructions are also ambiguous without
additional context. A practical system should allow users to inspect the
parsed layout, clarify uncertain constraints, and revise individual objects.

\section{Conclusion}

Structured scene records provide a compact interface between language and
visual generation. They make spatial requirements easier to inspect and can
support targeted editing. The example content in this template should be
replaced with the report's own methods, evidence, and limitations.

\begin{thebibliography}{9}
\bibitem[Radford et~al.(2021)]{radford2021clip}
  Alec Radford et~al.
  \newblock Learning transferable visual models from natural language supervision.
  \newblock In \emph{Proceedings of the International Conference on Machine
  Learning}, pages 8748--8763, 2021.
\end{thebibliography}

\ifbackauthorlist
\section*{Author List}
{\rmfamily\bfseries\ReportAuthors\par}
\vspace{3pt}
{\rmfamily\ReportInstitution\par}
\vspace{3pt}
{\rmfamily\itshape\small\ReportAuthorNotes\par}
\vspace{3pt}
{\rmfamily\small\textbf{Correspondence:} \ReportContact\par}
\fi

\end{document}
